% ##############################################################################
\section{NoesisServer: fulfillment and monitoring}
\label{sec:noesis_server}

\NoesisServer{} is the layer that turns a matched contract
(Definition~\ref{def:contract}) into actual inference requests against
real providers, and turns the resulting traffic into three things: a
measurement of whether the contract was honored, a signal for
difficulty-aware routing, and a corpus for downstream distillation.

% ==============================================================================
\subsection{Gateway architecture}
\label{sec:gateway_architecture}

\NoesisServer{} is, at its core, a lightweight LLM API gateway modeled on
multi-provider routers such as OpenRouter: it exposes a single API,
proxies each request to one or more underlying providers, and streams the
response back to the caller. Two properties distinguish it from a plain
pass-through router. First, every request is tagged with the contract
$\kappa$ it is served under, if any, so that fulfillment can be attributed
(Section~\ref{sec:fulfillment_monitoring}). Second, every prompt/response
pair is logged with its provider, model, measured latency, measured cost,
and, when available, a user feedback signal, regardless of whether the
request was placed under a contract. The storage schema is intentionally
minimal:

\[
  (\text{prompt}, \text{response}, \text{provider}, \text{model},
   \text{latency}, \text{cost}, \kappa\text{ or } \bot)
\]

This logging is passive by default: it does not change routing behavior
and exists purely to build a reusable dataset for evaluation and
distillation (Section~\ref{sec:distillation}). Privacy and data-handling
safeguards (PII scrubbing, opt-in/opt-out) are necessary before logging
real traffic and are discussed as an open question in
Section~\ref{sec:open_questions}.

% ==============================================================================
\subsection{Contract fulfillment monitoring}
\label{sec:fulfillment_monitoring}

For a request $x$ served under contract $\kappa$, let $\hat c(x) \in
\calC$ denote the capability tier actually delivered and $\hat\ell(x)$ the
actually measured latency, both estimated from the response (e.g., via a
verifier model, a benchmark probe, or agreement with a stronger reference
model). Request $x$ is a \emph{success} under $\kappa$ if

\begin{equation}
  \label{eq:success_indicator}
  s(x) = \mathbb{1}\!\left[
    \hat c(x) \succeq \clevelattr(\kappa)
    \;\wedge\;
    \hat\ell(x) \leq \lmaxattr(\kappa)
  \right].
\end{equation}

\begin{definition}[Measured reliability]
\label{def:measured_reliability}
For a fulfillment window $W$, the set of requests served under contract
$\kappa$, the measured reliability is
\begin{equation}
  \label{eq:measured_reliability}
  \hat r(\kappa, W) = \frac{1}{|W|} \sum_{x \in W} s(x).
\end{equation}
\end{definition}

A naive violation rule, flagging $\kappa$ whenever
$\hat r(\kappa, W) < \rminattr(\kappa)$, conflates genuine under-delivery
with sampling noise: $\hat r(\kappa, W)$ is a sample average of an
unobserved true success rate, and a small $|W|$ makes false accusations
likely. We instead treat fulfillment monitoring as a one-sided hypothesis
test: the null hypothesis $H_0$ is that the seller is compliant, i.e., its
true success rate is at least $\rminattr(\kappa)$, and $\kappa$ is flagged
as violated only if $H_0$ is rejected at a chosen confidence level
$1 - \delta$. A simple normal-approximation lower confidence bound gives

\begin{equation}
  \label{eq:reliability_lower_bound}
  \hat r_{\text{lower}}(\kappa, W)
  = \hat r(\kappa, W)
  - z_{1 - \delta}
    \sqrt{
      \frac{\hat r(\kappa, W) \, (1 - \hat r(\kappa, W))}{|W|}
    },
\end{equation}
where $z_{1 - \delta}$ is the corresponding standard normal quantile, and
$\kappa$ is flagged as violated when
$\hat r_{\text{lower}}(\kappa, W) < \rminattr(\kappa)$. For small $|W|$, an
exact interval (e.g., Clopper--Pearson) is preferable to the normal
approximation used here for exposition. Flagged violations are reported
to \NoesisMarket{} and feed the reputation update of
Equation~\eqref{eq:reputation_update}.

\begin{example}[Fulfillment reporting]
\label{ex:fulfillment_reporting}
A contract promising 99.9\% reliability at \texttt{frontier} capability is
served by three providers over a window of 5,000 requests. \NoesisServer{}
computes $\hat r(\kappa, W) = 0.994$ and, applying
Equation~\eqref{eq:reliability_lower_bound}, finds
$\hat r_{\text{lower}}(\kappa, W) < 0.999$: the shortfall is unlikely to be
sampling noise, so the contract is flagged as violated and reported back
to \NoesisMarket.
\end{example}

% ==============================================================================
\subsection{Difficulty-aware routing}
\label{sec:routing}

Not every request needs the strongest available model. \NoesisServer{}
trains a difficulty estimator
\begin{equation}
  \label{eq:difficulty_estimator}
  \hat d : \calP \to \calC
\end{equation}
mapping a prompt $x \in \calP$ to the cheapest capability tier expected to
answer it correctly, and routes according to
\begin{equation}
  \label{eq:routing_policy}
  \pi(x) = \min \{ c \in \calC : c \succeq \hat d(x) \}.
\end{equation}

Routing is only worthwhile if the estimator is cheaper than the savings it
produces. Let $\text{cost}(c)$ denote the per-request cost of serving tier
$c$, non-decreasing in $c$, let $c_K$ denote the strongest tier, and let
$c_{\hat d}$ denote the per-request cost of evaluating $\hat d$ itself.
Routing yields a net saving over an always-route-to-$c_K$ baseline exactly
when

\begin{equation}
  \label{eq:routing_saving_condition}
  c_{\hat d} + \EE_x [\, \text{cost}(\pi(x)) \,]
  <
  \text{cost}(c_K).
\end{equation}

Inequality~\eqref{eq:routing_saving_condition} makes explicit the trap
flagged in the underlying design notes: an estimator that is itself
expensive, or inaccurate enough to route hard requests to weak tiers and
force costly retries, can cost more than always calling the strong model.
Measuring the left-hand side against the right-hand side on real traffic
is a prerequisite for deploying routing in production, and is deferred to
the implementation plan in Section~\ref{sec:roadmap}.

% ==============================================================================
\subsection{Answer fusion}
\label{sec:fusion}

Fusion is the complementary lever to routing. Instead of picking one tier
per request, \NoesisServer{} may query a set of providers $S(x) \subseteq
\calM$ and combine their answers through an aggregator
\begin{equation}
  \label{eq:fusion_aggregator}
  \hat y(x) = g\big( \{ f_m(x) \}_{m \in S(x)} \big),
\end{equation}
where $g$ may be a majority vote, a verifier model that scores and selects
among candidates, or a learned combiner trained on logged outcomes.

Routing and fusion trade cost for accuracy in opposite directions relative
to a single-strong-model baseline: routing accepts a higher risk of a
lower-quality answer on some fraction of requests in exchange for a lower
expected cost, while fusion accepts a higher expected cost, from querying
multiple providers per request, in exchange for a lower risk of a
low-quality answer. Which lever is preferable for a given contract depends
on the contract's price versus its $\rminattr$ term: a low-price,
low-$\rminattr$ contract favors routing, while a high-$\rminattr$ contract
may only be economically served through fusion.

% ==============================================================================
\subsection{Training signal and distillation}
\label{sec:distillation}

The difficulty estimator of Equation~\eqref{eq:difficulty_estimator}
needs a training signal, which the logged traffic supplies for free.

\begin{definition}[Agreement label]
\label{def:agreement_label}
For a prompt $x$ answered by both a cheap model $f_{\text{cheap}}$ and a
strong model $f_{\text{strong}}$, the agreement label is
\begin{equation}
  \label{eq:agreement_label}
  y(x) = 1 - \mathbb{1}\big[ f_{\text{cheap}}(x) \equiv f_{\text{strong}}(x) \big],
\end{equation}
where $\equiv$ denotes exact match or a similarity threshold, used as a
proxy label for ``this prompt was hard'' ($y(x) = 1$) versus ``easy''
($y(x) = 0$).
\end{definition}

Definition~\ref{def:agreement_label} is a proxy, not a ground-truth
correctness label: the two models can agree on a wrong answer, or
disagree while the cheap model happens to be correct. A genuine quality
label requires independent verification, for instance from a pairwise
preference arena, which is a companion research direction discussed in
Section~\ref{sec:related_work}. With that caveat, the logged corpus
\[
  \calD = \{ (x_i, f_{\text{strong}}(x_i)) \}_{i=1}^N
\]
also supports knowledge distillation of a smaller model $f_\theta$ trained
to approximate $f_{\text{strong}}$'s outputs on the observed task
distribution \cite{hinton2015}, offering a second, independent way for
\NoesisServer{} to reduce cost on traffic it has already observed.
